% Copyright (c) 2026 Eugene Joseph M. Ragasa.
% SPDX-License-Identifier: Apache-2.0
% Author: Eugene Joseph M. Ragasa
% Formal proof completed: 2026-10-28
\documentclass[11pt]{article}
\usepackage[T1]{fontenc}\usepackage{amsmath,amssymb}\usepackage[margin=1in]{geometry}\usepackage[hidelinks]{hyperref}
\title{Proof Dossier: Pointwise Three-Channel Bound}\author{Eugene Joseph M. Ragasa}\date{Formal proof completed: 2026-10-28}
\begin{document}\raggedright\maketitle
\section{Identity and status}
ID \texttt{BG-HILBERT-THM-019}; status \texttt{proved}. Lean declaration:
\par\noindent{\tiny\path{ProjectKoios.BandGauge.Hilbert.norm_totalResidual_le_add}}.
\section{Statement and proof}
At each coordinate, $\lVert H-T\rVert\leq\lVert H-I\rVert+\lVert I-T\rVert$. The proof rewrites by the exact residual decomposition and applies the norm triangle inequality.
\section{Boundary and audit}
This is an upper bound, not equality of scalar error channels. Axioms are \texttt{propext}, \texttt{Classical.choice}, and \texttt{Quot.sound}; no admissions. Lean \texttt{4.35.0}, Mathlib {\scriptsize\path{56a005855ea17e0b0958f73ec801ae140b0a1679}}; AI-assisted, kernel-checked, no independent expert review recorded.
\end{document}
